
\documentclass[10pt]{article}
\usepackage{graphicx}
\usepackage{amsmath}
\usepackage{amssymb}
\usepackage{amsfonts}
\usepackage{textcomp}
\usepackage{amsthm}

%% ---- Theorem Environments ---------------------------------
\theoremstyle{plain}
\newtheorem{theorem}{Theorem}[section]
\newtheorem{lemma}[theorem]{Lemma}
\newtheorem{corollary}[theorem]{Corollary}
\newtheorem{proposition}[theorem]{Proposition}

\theoremstyle{definition}
\newtheorem{definition}[theorem]{Definition}
\newtheorem{example}[theorem]{Example}

\theoremstyle{remark}
\newtheorem{remark}[theorem]{Remark}

\title{Equations \& Theorems for Mathematics}
\author{A. E. Mahrouss\\amlal@nekernel.org}
\date{\today}

\begin{document}

\maketitle
\vspace{-0.5em}
\vspace{0.8em}
\rule{\linewidth}{0.4pt}
\vspace{1em}

\section{Equations \& Theorems}

\begin{definition}
	Let the following:
	\begin{align*}
		\int (\lambda + \mu)\sqrt{ax^2 + bx  +c} \operatorname{dx}.
	\end{align*}
	We define the three general equations factoring $V'$:
	\begin{align*}
		V' &= (\lambda + \mu)\sqrt{ax^2 + bx  +c}x \int x \operatorname{d}(\lambda + \mu)\sqrt{ax^2 + bx  +c};\\
		\frac{V'}{(\lambda + \mu)} &= \sqrt{ax^2 + bx  +c}x \int x \operatorname{d}(\lambda + \mu)\sqrt{ax^2 + bx  +c};\\
		\sqrt{ax^2 + bx  +c}\frac{V'}{(\lambda + \mu)} &= (\sqrt{ax^2 + bx  +c})^{2}x \int x \operatorname{d}(\lambda + \mu)\sqrt{ax^2 + bx  +c}.
	\end{align*}
\end{definition}
\begin{definition}
	Let the following:
	\begin{align*}
		O' + (\lambda + \mu + x)^{n-2} &= (\lambda + \mu + x)^{n-1} + (\lambda + \mu + x)^{n};\\
		O' &= \frac{(\lambda + \mu + x)^{n-1} + (\lambda + \mu + x)^{n}}{(\lambda + \mu + x)^{n-2}}; O' \neq 0.
	\end{align*}
\end{definition}

\nocite{*}

\bibliographystyle{acm}
\bibliography{A_Course_of_Pure_Mathematics, General_Theory_Dirichlet_Series, Riemann}

\end{document}